\documentclass[runningheads]{llncs}

\usepackage[utf8]{inputenc}
\usepackage[T2A]{fontenc}
\usepackage[russian]{babel}
\usepackage{amsmath,amssymb}
\usepackage{graphicx}
\graphicspath{{figures/}}
\usepackage{multirow}
\usepackage{array}
\usepackage{hyperref}
\usepackage{algorithm}
\usepackage{algorithmic}
\usepackage{url}
\usepackage{microtype}
\usepackage{xcolor}
\usepackage{tikz}
\usetikzlibrary{arrows.meta,positioning,shapes.geometric,calc,fit,backgrounds}
\definecolor{modblue}{RGB}{200,220,240}
\definecolor{modgreen}{RGB}{200,235,200}
\definecolor{modorange}{RGB}{255,225,180}
\definecolor{modred}{RGB}{245,200,200}
\definecolor{modgray}{RGB}{235,235,235}
\definecolor{decyellow}{RGB}{255,245,180}

\renewcommand\UrlFont{\rmfamily}

\hyphenation{Graph-Code-BERT ex-plo-it ex-plo-it-шаб-ло-нов
sand-box вери-фи-ка-ции вери-фи-ка-ция ин-тер-фе-ри-ро-вать
клас-си-фи-ка-ции клас-си-фи-ка-ци-он-ной
экс-пе-ри-мен-таль-ных дис-кри-ми-на-тив-ных
пре-до-бу-чен-ной эмбед-дин-гов сани-тай-зе-ров
не-ста-биль-нос-ти execution-based
FAISS-ин-дек-се FAISS-по-ис-ком}

\begin{document}

\title{GNN-ReGVD+: конвейер обнаружения уязвимостей в коде C/C++ с execution-based верификацией}

\titlerunning{GNN-ReGVD+: конвейер обнаружения уязвимостей}

\author{А.\,В.~Антонов\inst{1} \and
В.\,С.~Буровин\inst{1} \and
В.\,С.~Воеводкин\inst{1}}

\authorrunning{А.\,В. Антонов et al.}

\institute{Национальный исследовательский университет <<Высшая школа экономики>>,\\
Нижегородский филиал, г.~Нижний Новгород, Россия\\
\email{\{avantonov,vsburovin\}@edu.hse.ru, b.voevodkin@gmail.com}}

\maketitle

% ====================================================================
%                        ABSTRACT
% ====================================================================
\begin{abstract}
Автоматическое обнаружение уязвимостей в исходном коде C/C++ на основе графовых нейронных сетей характеризуется значительным уровнем ложных срабатываний при отсутствии execution-based подтверждения эксплуатируемости. В данной работе предлагается конвейер GNN-ReGVD+, объединяющий четыре модуля: (1)~нейросетевую детекцию на основе модели ReGVD с параметро-эффективной адаптацией LoRA и гибридным скорингом через FAISS; (2)~извлечение релевантных exploit-шаблонов из векторной базы; (3)~каскадную адаптацию шаблонов с применением больших языковых моделей; (4)~execution-based верификацию в Docker-изолированной среде с мульти-санитайзерной поддержкой и coverage-guided фаззингом. Экспериментальная оценка на датасетах Devign (27\,318 функций) и MegaVul (339\,548 функций) показала, что гибридная архитектура не снижает базовые показатели модели ($\Delta < 0{,}5$~п.п.), LoRA-адаптация 0,12\% параметров обеспечивает прирост ROC-AUC на $+0{,}179$ при переносе на новый домен со снижением затрат VRAM в~1,8~раза, а модуль sandbox-верификации покрывает 18 типов уязвимостей (23~CWE) с четырёхуровневой классификацией доказательств. Система реализована в виде 9\,372 строк кода Python с 203 автоматическими тестами.

\keywords{обнаружение уязвимостей \and графовые нейронные сети \and LoRA \and FAISS \and sandbox-верификация \and фаззинг \and DevSecOps}
\end{abstract}

% ====================================================================
%                     1. ВВЕДЕНИЕ
% ====================================================================
\section{Введение}

Рост объёмов программного обеспечения и частоты кибератак обостряет потребность в средствах автоматического обнаружения уязвимостей в исходном коде~\cite{nist800}. По данным OWASP~\cite{owasp2024}, уязвимости в коде остаются одним из ведущих векторов атак на веб-приложения и системную инфраструктуру, при этом число публикуемых CVE ежегодно возрастает. Традиционные инструменты статического анализа (SAST), такие как SonarQube, CodeQL, Coverity и Semgrep, обеспечивают интерпретируемость и воспроизводимость результатов, однако характеризуются высоким уровнем ложных срабатываний~--- от 30 до 80\%~\cite{sultana2024}~--- и зависимостью от вручную составленных правил. Каждое правило требует экспертного знания конкретного класса уязвимостей, что делает масштабирование инструментов трудоёмким и подверженным ошибкам.

Методы на основе глубокого обучения представляют альтернативный подход к этой задаче. Графовые нейронные сети (GNN)~\cite{nguyen2022regvd,zhou2019devign} моделируют структурные зависимости программного кода, предобученные модели кода (GraphCodeBERT~\cite{guo2021graphcodebert}, CodeLlama, DeepSeek-Coder) улавливают семантические паттерны из больших корпусов~\cite{ding2024}. Эти подходы повышают точность детекции по сравнению с ручными правилами, однако разделяют общее фундаментальное ограничение: все они производят лишь бинарную классификацию фрагмента кода (уязвим / не уязвим), не предоставляя execution-based доказательств эксплуатируемости. Без таких доказательств разработчик вынужден вручную анализировать каждое предупреждение, что снижает доверие к автоматизированным средствам и замедляет процесс устранения дефектов~\cite{gajjar2026vbf}. Как следствие, значительная часть предупреждений игнорируется, а реальные уязвимости могут быть пропущены среди потока ложных срабатываний.

Данная работа направлена на преодоление указанного ограничения. Мы предлагаем конвейер GNN-ReGVD+, в котором нейросетевая детекция дополняется автоматической генерацией proof-of-concept эксплойтов и их верификацией в изолированной среде. Конвейер построен по принципу воронки: лёгковесная GNN-детекция быстро отсеивает безопасный код, а ресурсоёмкие операции (LLM-адаптация эксплойтов, Docker-верификация) применяются только к подозрительным фрагментам. Такая архитектура обеспечивает баланс между скоростью анализа и глубиной верификации.

Основные вклады работы:

\begin{enumerate}
    \item Интегрированная четырёхмодульная архитектура, объединяющая GNN-детекцию с гибридным скорингом, FAISS-поиск exploit-шаблонов, каскадную LLM-адаптацию и sandbox-верификацию с параллельным LLM-контуром триажа.
    \item Гибридный метод скоринга ReGVD+FAISS с LoRA-адаптацией (0,12\% параметров), обеспечивающей прирост ROC-AUC на $+0{,}179$ при domain transfer и снижение затрат VRAM в~1,8~раза.
    \item Мульти-санитайзерная execution-based верификация с покрытием 18~типов уязвимостей через три профиля (ASan+UBSan, MSan, TSan), coverage-guided фаззинг и четырёхуровневую классификацию доказательств с автоматическим CWE-маппингом (23~записи).
    \item Воспроизводимая реализация~--- 9\,372 строки кода Python, 203~автоматических теста, CLI-интерфейс для интеграции в CI/CD-инфраструктуру.
\end{enumerate}

Статья организована следующим образом. В разделе~2 рассматриваются существующие подходы к детекции уязвимостей и связанные работы. В разделе~3 описывается архитектура конвейера GNN-ReGVD+. В разделе~4 приводится описание реализации и тестирования. В разделе~5 представлены результаты экспериментальной оценки. В разделе~6 обсуждаются полученные результаты и ограничения. Раздел~7 содержит заключение.


% ====================================================================
%                  2. ОБЗОР СУЩЕСТВУЮЩИХ ПОДХОДОВ
% ====================================================================
\section{Обзор существующих подходов}

Задача автоматического обнаружения уязвимостей в исходном коде имеет долгую историю исследований. В данном разделе мы рассматриваем эволюцию подходов от статического анализа до современных методов на основе глубокого обучения, а также формулируем исследовательский разрыв, который мотивирует данную работу.

\textbf{Статический анализ и классическое машинное обучение.}
Ранние подходы к автоматическому обнаружению уязвимостей основывались на ручном написании правил и паттернов. Инструменты SAST (SonarQube, Coverity, CodeQL, Semgrep) применяют фиксированные правила к абстрактному синтаксическому дереву или потоку данных программы. Несмотря на высокую интерпретируемость, они ограничены теми классами уязвимостей, для которых составлены правила, и генерируют значительное число ложных срабатываний~\cite{sultana2024}. Классические методы машинного обучения (Random~Forest, SVM, логистическая регрессия) на основе инженерных признаков достигали лишь ROC-AUC порядка 0,46--0,50, не улавливая семантические закономерности уязвимого кода. Переход к глубокому обучению на последовательностях (LSTM, TextCNN, Transformer) позволил отказаться от ручного конструирования признаков, однако линейная обработка токенов не отражала структурные зависимости программ~--- управляющий поток, зависимости по данным, вложенность вызовов.

\textbf{Графовые нейронные сети.}
Графовые нейронные сети (GNN) стали ответом на это ограничение, представляя программу в виде графа, узлы которого соответствуют элементам кода, а рёбра~--- структурным и семантическим связям между ними. Модель Devign~\cite{zhou2019devign} продемонстрировала эффективность такого подхода, комбинируя представления графа управляющего потока (CFG) и графа зависимостей по данным (PDG) с gated GNN-слоями. ReGVD~\cite{nguyen2022regvd} упростила подход, предложив построение графа непосредственно из токенов кода через механизм скользящего окна фиксированного размера, что устраняет необходимость в отдельном этапе синтаксического анализа. Модель использует residual connections для стабилизации обучения глубоких графовых сетей и soft attention pooling для агрегации узловых представлений, достигнув F1~$\approx 0{,}72$ на датасете Devign.

Работы по специализированной детекции, в частности обнаружение уязвимостей use-after-free через deep graph learning~\cite{chen2024uaf}, показали перспективность GNN для отдельных классов дефектов. Вместе с тем систематическое исследование Ding и др.~\cite{ding2024} выявило, что при корректном дизайне экспериментов (устранение дублирования, контроль распределения данных) заявленные метрики многих моделей существенно снижаются, что указывает на проблемы воспроизводимости в данной области. Несмотря на это, GNN остаются одним из наиболее перспективных направлений благодаря способности моделировать структурные свойства программ.

\textbf{Предобученные языковые модели кода.}
Предобученные модели (GraphCodeBERT~\cite{guo2021graphcodebert}, CodeLlama, DeepSeek-Coder, StarCoder) достигают ROC-AUC 80--83\% в zero-shot режиме~\cite{ding2024}, что делает их привлекательными для детекции уязвимостей без обучения на специализированных данных. SAST-Genius~\cite{agrawal2025sast} применяет большие языковые модели (LLM) для триажа алертов статического анализа CodeQL, извлекая графовый контекст для повышения точности классификации. Нейро-символический подход Li~и~др.~\cite{li2025neurosymbolic} комбинирует LLM-паттерны уязвимостей со статическим анализом. Сравнительные исследования~\cite{sultana2024,makhyari2024} подтвердили, что LLM обнаруживают семантически сложные паттерны, но страдают от нестабильности вывода и высокой вычислительной стоимости. Все рассмотренные подходы ограничены бинарной классификацией кода без подтверждения эксплуатируемости.

\textbf{Execution-based верификация и фаззинг.}
Концепция Verify Before You Fix~\cite{gajjar2026vbf} формализует идею execution-based валидации через Docker-sandbox с классификацией доказательств (CONFIRMED / SUGGESTIVE / INCONCLUSIVE). Данный подход предлагает сначала подтвердить наличие уязвимости путём запуска exploit-кода в изолированной среде и лишь затем приступать к исправлению. Фаззинг эволюционировал от случайной мутации входных данных (AFL) к coverage-guided подходам (LibFuzzer, Honggfuzz), которые используют информацию о покрытии кода для направления генерации тестов. Санитайзеры памяти (ASan, MSan, TSan) дополняют фаззинг, обнаруживая ошибки, которые не приводят к немедленному аварийному завершению, но представляют серьёзные уязвимости~--- обращение к неинициализированной памяти, гонки данных, выход за границы буфера без SIGSEGV.

\textbf{Параметро-эффективная адаптация.}
LoRA (Low-Rank Adaptation)~\cite{hu2021lora} добавляет компактные обучаемые матрицы низкого ранга ($r \ll \min(d, k)$) к attention-слоям предобученной модели, сохраняя остальные параметры замороженными. В отличие от полного fine-tuning, LoRA требует хранения и обновления лишь малой доли параметров, что значительно снижает требования к вычислительным ресурсам. В контексте детекции уязвимостей LoRA перспективна для domain transfer: модель, обученная на одном наборе данных, адаптируется к новой кодовой базе с иным распределением уязвимостей без полного переобучения.

\textbf{Исследовательский разрыв.}
Сравнение существующих парадигм (табл.~\ref{tab:paradigms}) выявляет ключевой разрыв: ни один из рассмотренных подходов не объединяет быструю GNN-детекцию с execution-based верификацией. Классический ML ограничен точностью, GNN обеспечивают хороший баланс скорости и качества, LLM~--- наилучшую точность, но при высоких вычислительных затратах. При этом все три парадигмы завершают работу на этапе классификации, не предоставляя доказательств эксплуатируемости. GNN-ReGVD+ заполняет этот разрыв, реализуя полный цикл от детекции до sandbox-верификации с автоматическим CWE-маппингом.

\begin{table}
\centering
\caption{Сравнение парадигм автоматического анализа уязвимостей}
\label{tab:paradigms}
\small
\begin{tabular}{|p{2.0cm}|p{2.6cm}|p{2.8cm}|p{3.0cm}|}
\hline
\textbf{Критерий} & \textbf{Классический ML} & \textbf{GNN (ReGVD)} & \textbf{LLM (CodeLlama)} \\
\hline
Представление & Инженерные признаки & Графы (sliding window) & Последовательности токенов \\
Точность & ROC-AUC $\approx$ 0,46--0,50 & Acc $\approx$ 62--64\% & ROC-AUC 80--83\% \\
Латентность & Микросекунды & Миллисекунды & Секунды \\
Адаптивность & Низкая & Средняя (LoRA) & Высокая (zero-shot) \\
Верификация & Нет & Нет & Нет \\
\hline
\end{tabular}
\end{table}


% ====================================================================
%              3. АРХИТЕКТУРА КОНВЕЙЕРА
% ====================================================================
\section{Архитектура конвейера GNN-ReGVD+}

В данном разделе описывается архитектура предлагаемого конвейера. Система состоит из четырёх последовательных модулей: GNN-детекция с гибридным скорингом, извлечение exploit-шаблонов из векторной базы, каскадная LLM-адаптация шаблонов и sandbox-верификация в Docker-изолированной среде. Параллельно с основным конвейером функционирует LLM-голова триажа, обеспечивающая независимый канал семантического анализа. Общая архитектура представлена на рис.~\ref{fig:pipeline}.

\begin{figure}
\centering
\resizebox{\textwidth}{!}{%
\begin{tikzpicture}[
    module/.style={rectangle, draw=black!60, rounded corners=4pt, minimum height=1.2cm,
        text width=3.2cm, align=center, font=\footnotesize, line width=0.8pt},
    arr/.style={-{Stealth[length=4pt]}, semithick},
    darr/.style={-{Stealth[length=4pt]}, semithick, dashed, draw=black!40},
    node distance=0.8cm
]
\node[module, fill=modblue] (m1) {\textbf{Module 1}\\GNN Detection\\[-1pt]{\scriptsize ReGVD + LoRA + FAISS}};
\node[module, fill=modblue, right=of m1] (m2) {\textbf{Module 2}\\Exploit Retrieval\\[-1pt]{\scriptsize FAISS $k$-NN, 20 tpl.}};
\node[module, fill=modorange, right=of m2] (m3) {\textbf{Module 3}\\LLM Adaptation\\[-1pt]{\scriptsize LLM $\to$ Templ.}\\[-1pt]{\scriptsize $\to$ Generic}};
\node[module, fill=modgreen, right=of m3] (m4) {\textbf{Module 4}\\Sandbox Verification\\[-1pt]{\scriptsize Docker + ASan /}\\[-1pt]{\scriptsize MSan / TSan + LibFuzzer}};

\node[left=0.6cm of m1, font=\footnotesize, align=center] (input) {Source\\Code\\(C/C++)};
\draw[arr] (input) -- (m1);
\draw[arr] (m1) -- node[above, font=\scriptsize] {emb.\ (512-d)} (m2);
\draw[arr] (m2) -- (m3);
\draw[arr] (m3) -- (m4);

\node[rectangle, draw=black!60, rounded corners=4pt, right=0.8cm of m4, align=center,
    font=\footnotesize, minimum height=1.0cm, text width=2.2cm, fill=modgray, line width=0.8pt]
    (verdict) {\textbf{Verdict}\\[-1pt]{\scriptsize CONFIRMED}\\[-1pt]{\scriptsize SUGGESTIVE}\\[-1pt]{\scriptsize NEUTRAL / SAFE}};
\draw[arr] (m4) -- (verdict);

\node[module, fill=modred, text width=2.6cm, above=0.8cm of m2] (triage)
    {\textbf{LLM Triage}\\[-1pt]{\scriptsize Claude Sonnet}\\[-1pt]{\scriptsize + CodeQL}};
\draw[arr] (m1.north) |- (triage.west);
\draw[darr] (triage.east) -| node[near start, above, font=\scriptsize, text=black!50]
    {cross-validation} (verdict.north);

\draw[darr] (m4.south) -- ++(0,-0.6) -| node[near start, below, font=\scriptsize,
    text=black!50] {compile error $\to$ retry} (m3.south);
\end{tikzpicture}%
}
\caption{Архитектура конвейера GNN-ReGVD+. Исходный код последовательно проходит через четыре модуля. Параллельно работает LLM-голова триажа. Пунктирная стрелка обозначает итеративное уточнение при ошибке компиляции harness}
\label{fig:pipeline}
\end{figure}

\subsection{Модуль детекции: ReGVD с LoRA-адаптацией и гибридным скорингом}

В основе модуля детекции лежит модель ReGVD~\cite{nguyen2022regvd}, расширенная двумя ключевыми компонентами: LoRA-адаптерами для параметро-эффективного domain transfer и FAISS-головой для гибридного скоринга с определением типа уязвимости (рис.~\ref{fig:model}).

\begin{figure}
\centering
\scalebox{0.82}{%
\begin{tikzpicture}[
    block/.style={rectangle, draw=black!60, rounded corners=3pt, minimum width=3.0cm,
        minimum height=0.55cm, align=center, font=\footnotesize, line width=0.8pt},
    sblock/.style={rectangle, draw=black!60, rounded corners=3pt, minimum width=2.2cm,
        minimum height=0.5cm, align=center, font=\scriptsize, line width=0.8pt},
    arr/.style={-{Stealth[length=3pt]}, semithick},
    node distance=0.35cm
]
\node[block, fill=modgray] (input) {Source Code Tokens};
\node[block, fill=modblue, below=of input] (tok) {BPE Tokenizer (GraphCodeBERT)};
\draw[arr] (input) -- (tok);
\node[block, fill=modblue, below=of tok] (graph) {Sliding Window Graph ($w{=}5$)};
\draw[arr] (tok) -- (graph);
\node[block, fill=modblue!70!black!20, below=0.4cm of graph, minimum width=3.6cm,
    minimum height=0.65cm, font=\footnotesize\bfseries] (gcb) {GraphCodeBERT (frozen, 125M)};
\draw[arr] (graph) -- (gcb);
\node[sblock, fill=modorange, right=0.2cm of gcb] (lora) {LoRA\\[-1pt]$r{=}8,\;\alpha{=}16$\\[-1pt]152.8K params};
\draw[semithick] (gcb.east) -- (lora.west);
\node[block, fill=modblue, below=0.4cm of gcb] (regcn) {ReGCN ($L{=}2$, $h{=}128$) + Residual};
\draw[arr] (gcb) -- (regcn);
\node[block, fill=modblue, below=of regcn] (pool) {Soft Attention Pooling};
\draw[arr] (regcn) -- (pool);

\node[sblock, fill=modgreen, below left=0.55cm and 0.8cm of pool] (cls) {Classification Head\\[-1pt]Linear $\to$ Sigmoid};
\node[sblock, fill=modorange, below right=0.55cm and 0.8cm of pool] (emb) {Embedding Head\\[-1pt]Linear $768{\to}512$, $L_2$};
\draw[arr] (pool.south) -- ++(0,-0.18) -| (cls.north);
\draw[arr] (pool.south) -- ++(0,-0.18) -| (emb.north);

\node[below=0.25cm of cls, font=\footnotesize\bfseries] (pcls) {$P_{\text{cls}}$};
\draw[arr] (cls) -- (pcls);

\node[sblock, fill=modorange, below=0.25cm of emb] (faiss) {FAISS IndexFlatIP\\[-1pt]$k$-NN search};
\draw[arr] (emb) -- (faiss);
\node[below=0.25cm of faiss, font=\footnotesize\bfseries, align=center] (sfaiss) {$S_{\text{FAISS}}$, vuln.\ type};
\draw[arr] (faiss) -- (sfaiss);

\node[block, fill=modgreen!60!black!10, below=3.2cm of pool, minimum width=4.8cm]
    (hybrid) {$S_{\text{hybrid}} = \beta \cdot P_{\text{cls}} + (1{-}\beta)\cdot S_{\text{FAISS}}$};
\draw[arr] (pcls.south) |- (hybrid.west);
\draw[arr] (sfaiss.south) |- (hybrid.east);
\end{tikzpicture}%
}
\caption{Архитектура модели GNNReGVD с двойным скорингом. GraphCodeBERT (frozen) с LoRA-адаптерами формирует $P_{\text{cls}}$; ReGCN + embedding head + FAISS формируют $S_{\text{FAISS}}$ и определяют тип уязвимости через majority vote}
\label{fig:model}
\end{figure}

\textbf{Базовая модель ReGVD.} Исходный код функции токенизируется с помощью BPE-токенизатора GraphCodeBERT~\cite{guo2021graphcodebert}. Последовательность токенов преобразуется в граф через механизм скользящего окна размера $w$: каждый токен $t_i$ соединяется с $w$ последующими токенами $t_{i+1}, \ldots, t_{i+w}$. Полученный граф обрабатывается двухслойной ReGCN ($L{=}2$, $h{=}128$) с residual connections. Soft attention pooling агрегирует узловые представления в единый вектор графа, который подаётся на классификационную голову.

\textbf{LoRA-адаптация.} Веса модели GraphCodeBERT (${\sim}125$M параметров) заморожены. К проекциям query и value во всех 12~attention-слоях добавлены LoRA-адаптеры ($r = 8$, $\alpha = 16$):
\begin{equation}
    \mathbf{h} = \mathbf{W}\mathbf{x} + \mathbf{B}(\mathbf{A}\mathbf{x}) \cdot \frac{\alpha}{r},
    \label{eq:lora}
\end{equation}
где $\mathbf{A} \in \mathbb{R}^{r \times k}$, $\mathbf{B} \in \mathbb{R}^{d \times r}$~--- обучаемые матрицы низкого ранга, $\mathbf{W} \in \mathbb{R}^{d \times k}$~--- замороженные веса исходной модели. Общее число обучаемых параметров составляет ${\sim}152{,}8$K (0,12\% от полной модели), что снижает требования к VRAM с ${\sim}14{,}2$~ГБ до ${\sim}7{,}8$~ГБ при дообучении и энергозатраты приблизительно в~8 раз. Данное сокращение ресурсов переводит дообучение из категории серверных GPU (RTX~3090 с 24~ГБ) в категорию потребительских GPU с 8~ГБ видеопамяти.

\textbf{FAISS-голова и гибридный скоринг.} Для каждой классифицированной функции embedding head (Linear 768$\to$512, $L_2$-нормализация) проецирует выход ReGCN в пространство для $k$-NN поиска с использованием библиотеки FAISS~\cite{johnson2019faiss}. Индекс типа \texttt{IndexFlatIP} обеспечивает точный поиск по внутреннему произведению нормализованных векторов. Гибридный скор объединяет вероятность классификатора и долю уязвимых функций среди $K$ ближайших соседей:
\begin{equation}
    S_{\text{hybrid}} = \beta \cdot P_{\text{cls}} + (1 - \beta) \cdot S_{\text{FAISS}}, \quad S_{\text{FAISS}} = \frac{1}{K}\sum_{j=1}^{K} \mathbb{I}(y_j = 1),
    \label{eq:hybrid}
\end{equation}
где $\beta$ подбирается через grid search на валидационной выборке. Тип уязвимости определяется через majority vote по метаданным $K$ ближайших уязвимых соседей из обучающей выборки, что позволяет классифицировать дефект (buffer~overflow, use-after-free, integer~overflow и~т.\,д.) без обучения отдельного мультиклассового классификатора.

Обучение модели выполняется с комбинированной функцией потерь, объединяющей бинарную кросс-энтропию для задачи классификации и контрастивную компоненту для формирования дискриминативных эмбеддингов:
\begin{equation}
    \mathcal{L}_{\text{total}} = \alpha \cdot \mathcal{L}_{\text{BCE}} + (1 - \alpha) \cdot \mathcal{L}_{\text{SupCon}},
    \label{eq:loss}
\end{equation}
где $\mathcal{L}_{\text{SupCon}}$~--- supervised contrastive loss~\cite{khosla2020supcon}, притягивающий эмбеддинги функций одного класса и отталкивающий эмбеддинги разных классов. Весовой коэффициент $\alpha = 0{,}7$ определён экспериментально.

\textbf{Двухконтурное обновление.} Система поддерживает два режима обновления. FAISS-индекс поддерживает hot-update (${\sim}0{,}18$~с на 100~эмбеддингов) для оперативного добавления новых CVE без переобучения модели. LoRA-адаптеры обеспечивают периодическое дообучение при накоплении достаточного объёма новых данных. Для предотвращения catastrophic forgetting при LoRA-дообучении используется replay-стратегия: 20\% каждого batch составляют примеры из исходного домена (Devign).


\subsection{Извлечение, адаптация и sandbox-верификация exploit-шаблонов}

Второй, третий и четвёртый модули конвейера обеспечивают полный цикл от обнаружения потенциальной уязвимости до execution-based подтверждения её эксплуатируемости.

\textbf{Извлечение exploit-шаблонов.} 512-мерный эмбеддинг из модуля детекции используется для поиска в отдельном FAISS-индексе exploit-шаблонов. Поиск выполняется в два этапа: сначала broad similarity ($5 \times k$ кандидатов), затем фильтрация по типу уязвимости и отбор top-$k$ (по умолчанию~3). База содержит 20~exploit-шаблонов в 10~категориях уязвимостей (табл.~\ref{tab:templates}): buffer overflow (stack и heap), use-after-free, format string, integer overflow, null pointer dereference, off-by-one, division by zero, uninitialized read и race condition. Каждый шаблон содержит параметризованный C-код с маркерами подстановки (\texttt{\{vulnerable\_code\}}, \texttt{\{func\_name\}}) и метаданные (тип уязвимости, CWE, описание механизма эксплуатации).

\begin{table}
\centering
\caption{Exploit-шаблоны по категориям уязвимостей}
\label{tab:templates}
\small
\begin{tabular}{|l|l|c|}
\hline
\textbf{Категория} & \textbf{Покрытие} & \textbf{Кол-во} \\
\hline
buffer\_overflow & strcpy, sprintf, memcpy & 3 \\
heap\_buffer\_overflow & malloc, realloc & 2 \\
use\_after\_free & basic UAF, double-free & 2 \\
format\_string & read (\%x), write (\%n) & 2 \\
integer\_overflow & signed, alloc & 2 \\
null\_pointer\_deref & basic, return-NULL & 2 \\
off\_by\_one & loop, strlen & 2 \\
division\_by\_zero & basic, modulo & 2 \\
uninitialized\_read & heap, stack & 2 \\
race\_condition & data race & 1 \\
\hline
\textbf{Итого} & & \textbf{20} \\
\hline
\end{tabular}
\end{table}

\textbf{Каскадная LLM-адаптация.} Адаптация exploit-шаблонов под конкретный код реализуется через каскад из трёх стратегий с fallback-логикой (рис.~\ref{fig:adaptation}).

\begin{figure}
\centering
\scalebox{0.82}{%
\begin{tikzpicture}[
    block/.style={rectangle, draw=black!60, rounded corners=3pt, minimum width=3.4cm,
        minimum height=0.55cm, align=center, font=\footnotesize, line width=0.8pt},
    decision/.style={diamond, draw=black!60, aspect=2.5, inner sep=1pt,
        font=\footnotesize, fill=decyellow, line width=0.8pt},
    success/.style={rectangle, draw=black!60, rounded corners=5pt, fill=modgreen,
        minimum width=1.6cm, minimum height=0.4cm, align=center, font=\footnotesize, line width=0.8pt},
    arr/.style={-{Stealth[length=3pt]}, semithick},
    darr/.style={-{Stealth[length=3pt]}, semithick, dashed, draw=black!40},
    node distance=0.35cm
]
\node[block, fill=modgray, rounded corners=6pt] (start) {Vulnerable code + exploit template};

\node[block, fill=modblue, below=0.45cm of start] (l1) {\textbf{Level 1:} LLM Adaptation\\[-1pt]{\scriptsize Qwen2.5-Coder / OpenAI API}};
\draw[arr] (start) -- (l1);

\node[decision, below=0.4cm of l1] (d1) {\scriptsize Compiles?};
\draw[arr] (l1) -- (d1);

\node[success, right=1.8cm of d1] (ok1) {Harness};
\draw[arr] (d1) -- node[above, font=\scriptsize] {yes} (ok1);

\draw[darr] (d1.west) -- ++(-1.4,0) node[left, font=\scriptsize, text=black!50]
    {retry ${\leq}\,3$} |- (l1.west);

\node[block, fill=modorange, below=0.6cm of d1] (l2) {\textbf{Level 2:} Template Substitution\\[-1pt]{\scriptsize \{vuln\_code\}, \{func\_name\}}};
\draw[arr] (d1) -- node[right, font=\scriptsize] {no} (l2);

\node[decision, below=0.4cm of l2] (d2) {\scriptsize Valid?};
\draw[arr] (l2) -- (d2);

\node[success, right=1.8cm of d2] (ok2) {Harness};
\draw[arr] (d2) -- node[above, font=\scriptsize] {yes} (ok2);

\node[block, fill=modred, below=0.6cm of d2] (l3) {\textbf{Level 3:} Generic Harness\\[-1pt]{\scriptsize 15 generators by vuln.\ type}};
\draw[arr] (d2) -- node[right, font=\scriptsize] {no} (l3);

\node[block, fill=modgreen!60!black!10, below=0.5cm of l3, rounded corners=6pt]
    (out) {$\to$ Sandbox Verification (Module 4)};
\draw[arr] (l3) -- (out);
\coordinate (routeR) at ([xshift=0.6cm]ok1.east);
\draw[arr] (ok1.east) -- (routeR) |- (out.east);
\draw[arr] (ok2.east) -- ([xshift=0.6cm]ok2.east) |- ([yshift=3pt]out.east);
\end{tikzpicture}%
}
\caption{Каскад адаптации exploit-шаблонов. При неудаче каждого уровня управление передаётся следующему: LLM-адаптация с итеративным уточнением $\to$ template-based подстановка $\to$ generic harness}
\label{fig:adaptation}
\end{figure}

\textit{Уровень~1: LLM-адаптация} с итеративным уточнением~\cite{gajjar2026vbf}. LLM (Qwen2.5-Coder-7B или OpenAI-compatible API) получает структурированный промпт, содержащий уязвимый код, тип уязвимости, найденный exploit-шаблон, а также taint-контекст (source$\to$sink пути), если он доступен. Для 14~типов уязвимостей определены специализированные подсказки (hints), направляющие модель на генерацию конкретного exploit-механизма. При ошибке компиляции сгенерированного harness система включает сообщение компилятора в следующий промпт и повторяет запрос (до~3~попыток), что позволяет LLM скорректировать код.

\textit{Уровень~2: Template-based подстановка.} Если LLM-адаптация не удалась (все попытки завершились ошибкой или LLM недоступна), выполняется механическая подстановка переменных в найденный exploit-шаблон: уязвимый код, имя функции, аргументы.

\textit{Уровень~3: Generic harness.} В случае неудачи обоих предыдущих уровней один из 15~специализированных генераторов harness (по одному на каждый поддерживаемый тип уязвимости) гарантированно возвращает минимальный валидный harness-код, вызывающий целевую функцию с потенциально опасными аргументами.

Метод \texttt{adapt\_multi()} генерирует несколько вариантов harness с дедупликацией по MD5-хешу, увеличивая вероятность успешной верификации.

\textbf{Sandbox-верификация.} Каждый сгенерированный harness компилируется и запускается в Docker-изолированной среде с жёсткими ограничениями безопасности: \texttt{network=none}, \texttt{cap-drop=ALL}, seccomp whitelist из 70+ разрешённых системных вызовов, ограничение памяти 1~ГБ и PID~limit~128. Такая конфигурация предотвращает выход вредоносного кода за пределы контейнера.

Система поддерживает три взаимоисключающих профиля санитайзеров (табл.~\ref{tab:sanitizers}) с автоматическим выбором на основе типа уязвимости. Профиль \texttt{asan\_ubsan} (ASan~+ UBSan) детектирует buffer overflow, use-after-free, null~pointer dereference и другие ошибки работы с памятью. Профиль \texttt{msan} (MemorySanitizer) специализирован на обнаружении чтения неинициализированной памяти (CWE-457). Профиль \texttt{tsan} (ThreadSanitizer) обнаруживает гонки данных (CWE-362) и взаимные блокировки (CWE-833). Взаимоисключаемость обусловлена несовместимостью санитайзеров на уровне инструментирования компилятора.

\begin{table}
\centering
\caption{Профили санитайзеров для sandbox-верификации}
\label{tab:sanitizers}
\small
\begin{tabular}{|l|p{3.2cm}|p{4.0cm}|}
\hline
\textbf{Профиль} & \textbf{Флаги} & \textbf{Детектирует} \\
\hline
asan\_ubsan & \texttt{-fsanitize= address, undefined} & buffer overflow, UAF, null deref, int overflow \\
msan & \texttt{-fsanitize= memory} & uninit. read (CWE-457) \\
tsan & \texttt{-fsanitize= thread} & data race (CWE-362), deadlock (CWE-833) \\
\hline
\end{tabular}
\end{table}

Верификация является двухфазной. На первой фазе harness запускается с заданными входными данными (статический payload) с timeout~1~с. Если санитайзер не обнаружил ошибку, запускается вторая фаза~--- coverage-guided фаззинг через LibFuzzer. Для этого выполняется автоматическая трансформация \texttt{main()} $\to$ \texttt{LLVMFuzzerTestOneInput()}: система анализирует сигнатуру \texttt{main()} и генерирует обёртку, передающую фаззированные данные в качестве аргументов целевой функции. LibFuzzer использует информацию о покрытии кода для направления мутаций, что увеличивает вероятность обнаружения уязвимости по сравнению со случайным тестированием.

\textbf{Классификация доказательств и CWE-маппинг.} Вывод санитайзеров классифицируется через 71~регулярный паттерн (ASan:~17, UBSan:~14, MSan:~5, TSan:~7, Crash:~16, Valgrind:~12) в четырёхуровневую шкалу с confidence scoring (табл.~\ref{tab:confidence}). Каждому паттерну присвоен базовый уровень confidence, отражающий надёжность обнаружения: санитайзеры AddressSanitizer и MemorySanitizer дают наиболее достоверные результаты (0,93--0,95), тогда как косвенные признаки (non-zero exit code) имеют значительно более низкий confidence (0,30).

Автоматический CWE-маппинг (23~записи) сопоставляет каждый обнаруженный паттерн санитайзера с идентификатором CWE. Например, паттерн \texttt{heap-buffer-overflow} маппируется на CWE-122, \texttt{use-after-poison}~--- на CWE-416, \texttt{data race}~--- на CWE-362.

Финальный вердикт для каждой проверяемой функции определяется иерархически: \textbf{CONFIRMED}~--- хотя бы один harness подтверждён санитайзером или crash-сигналом; \textbf{SUGGESTIVE}~--- обнаружены косвенные признаки (Valgrind leak, non-zero exit); \textbf{NEUTRAL}~--- sandbox-верификация не дала результатов; \textbf{SAFE}~--- $S_{\text{hybrid}} < \tau$ (функция не помечена как подозрительная).

\begin{table}
\centering
\caption{Confidence scoring по источнику доказательства}
\label{tab:confidence}
\small
\begin{tabular}{|l|c|c|}
\hline
\textbf{Источник} & \textbf{Base confidence} & \textbf{Уровень} \\
\hline
ASan (17 паттернов) & 0,95 & CONFIRMED \\
MSan (5 паттернов) & 0,93 & CONFIRMED \\
TSan (7 паттернов) & 0,91 & CONFIRMED \\
UBSan (14 паттернов) & 0,90 & CONFIRMED \\
Crash/Signal (16 паттернов) & 0,85 & CONFIRMED \\
Valgrind leak & 0,60 & SUGGESTIVE \\
Non-zero exit & 0,30 & SUGGESTIVE \\
Clean exit & 1,00 & SAFE \\
\hline
\end{tabular}
\end{table}

\textbf{LLM-триаж: параллельная голова верификации.}
Параллельно с sandbox-верификацией в конвейере функционирует вторая голова триажа на основе LLM (Claude Sonnet 4.6), работающая как независимый канал валидации. Для каждого подозрительного фрагмента система извлекает граф зависимостей через CodeQL, формирует структурированный промпт, содержащий исходный код, графовый контекст (потоки данных, вызовы функций) и предполагаемый тип уязвимости, и запрашивает LLM-классификацию (vulnerable~/~safe~/~uncertain).

Система реализует асимметричную двухголовую верификацию: sandbox-голова предоставляет execution-based evidence через санитайзеры и фаззинг, LLM-голова~--- семантический анализ с графовым контекстом. Перекрёстное подтверждение (обе головы классифицируют функцию как уязвимую) существенно повышает confidence итогового вердикта. Расхождение результатов (одна голова считает код уязвимым, другая~--- безопасным) сигнализирует о необходимости ручной проверки. LLM-голова особенно полезна в случаях, когда sandbox-верификация затруднена~--- например, при невозможности собрать harness из-за сложных зависимостей,~--- обеспечивая fallback-канал триажа.


% ====================================================================
%              4. РЕАЛИЗАЦИЯ И ТЕСТИРОВАНИЕ
% ====================================================================
\section{Реализация и тестирование}

Система реализована в виде модульной кодовой базы на языке Python\footnote{Исходный код: \url{https://github.com/deccersw/GNN-ReGVD-Plus}} объёмом ${\sim}9\,372$ строки с lazy initialization модулей (табл.~\ref{tab:codebase}). Архитектура следует принципу разделения ответственности: каждый модуль конвейера инкапсулирован в отдельном пакете с чётко определённым интерфейсом. Lazy initialization обеспечивает загрузку тяжёлых компонентов (модель GNN, FAISS-индексы, Docker-клиент) только при первом обращении, что ускоряет запуск и снижает потребление памяти при работе в sandbox-only режиме.

\begin{table}
\centering
\caption{Структура кодовой базы GNN-ReGVD+}
\label{tab:codebase}
\small
\begin{tabular}{|l|l|r|}
\hline
\textbf{Модуль} & \textbf{Основные компоненты} & \textbf{Строк} \\
\hline
\texttt{code/} & GNN-модель, FAISS-индекс, inference engine & ${\sim}2\,800$ \\
\texttt{exploit\_db/} & Exploit index, adapter, шаблоны & ${\sim}1\,500$ \\
\texttt{sandbox/} & Executor, evidence classifier, Docker & ${\sim}1\,500$ \\
\texttt{analysis/} & Taint analysis (Joern CPG / regex) & ${\sim}400$ \\
\texttt{scanner/} & Pipeline orchestrator, config, report & ${\sim}600$ \\
\texttt{tests/} & 203 автоматических теста & ${\sim}2\,500$ \\
CLI + evaluate & \texttt{scan\_cli.py}, \texttt{evaluate.py} & ${\sim}570$ \\
\hline
\textbf{Итого} & & ${\sim}$\textbf{9\,372} \\
\hline
\end{tabular}
\end{table}

CLI-интерфейс (\texttt{scan\_cli.py}) поддерживает три режима работы: анализ одиночной функции (интерактивный ввод или файл), batch-сканирование JSONL-датасетов и sandbox-only режим (верификация без GNN-детекции). Система автоматически определяет наличие GPU и переключается между CUDA- и CPU-инференсом. Конфигурация конвейера (пороги, гиперпараметры, пути к моделям) управляется через единый YAML-файл с валидацией схемы.

\textbf{Тестирование.} 203~автоматических теста покрывают все модули системы: sandbox (61~тест), enhancements (55), exploit adapter (36), pipeline (25), evaluation (18), FAISS (11). Тесты спроектированы для работы без GPU, Docker или обученной модели~--- все внешние зависимости замокированы. Время выполнения полного набора тестов составляет ${\sim}6$~секунд. Тестирование включает юнит-тесты отдельных компонентов (парсинг вывода санитайзеров, CWE-маппинг, генерация harness), интеграционные тесты (прохождение функции через полный конвейер с мокированными внешними сервисами) и regression-тесты для воспроизводимости результатов детекции.


% ====================================================================
%              5. ЭКСПЕРИМЕНТАЛЬНАЯ ОЦЕНКА
% ====================================================================
\section{Экспериментальная оценка}

\subsection{Протокол эксперимента и данные}

Экспериментальная программа включает четыре последовательных фазы, направленных на оценку каждого компонента гибридной архитектуры:

\begin{itemize}
    \item \textbf{Фаза~1:} Воспроизведение ReGVD baseline на датасете Devign для подтверждения корректности реализации.
    \item \textbf{Фаза~2:} Обучение гибридной модели (ReGVD + LoRA + FAISS) на Devign для оценки совместимости компонентов.
    \item \textbf{Фаза~3:} Zero-shot перенос модели, обученной на Devign, на датасет MegaVul для количественной оценки domain gap.
    \item \textbf{Фаза~4:} LoRA-адаптация на MegaVul для оценки эффективности параметро-эффективного domain transfer.
\end{itemize}

Использованы два датасета:
\begin{itemize}
    \item \textbf{Devign}~\cite{zhou2019devign}: 27\,318 функций на C из четырёх крупных open-source проектов (QEMU, FFmpeg, Linux Kernel, Wireshark). Распределение классов приблизительно сбалансировано (${\sim}45{,}6$\% уязвимых функций), что упрощает обучение, но не отражает реальное распределение уязвимостей в production-коде.
    \item \textbf{MegaVul}~\cite{ni2024megavul}: 339\,548 функций C/C++ из 992~open-source проектов. Содержит ${\sim}5{,}1$\% уязвимых функций (169~типов CWE), что значительно ближе к реальному распределению. Масштаб и разнообразие делают MegaVul значительно более сложным бенчмарком.
\end{itemize}

\textbf{Аппаратная среда}: 1$\times$ NVIDIA RTX 3090 (24~ГБ VRAM), AMD Ryzen~9, 64~ГБ RAM, Ubuntu~22.04.

\textbf{Гиперпараметры модели}: block\_size~=~400 (максимальная длина последовательности токенов), window\_size~=~5 (размер скользящего окна для построения графа), ReGCN ($L{=}2$ слоя, $h{=}128$ скрытых единиц), batch\_size~=~128, оптимизатор AdamW ($\text{lr} = 5 \times 10^{-4}$, weight decay $= 10^{-2}$), весовой коэффициент контрастивной потери $\alpha = 0{,}7$, температура SupCon $\tau = 0{,}07$.

\textbf{Гиперпараметры LoRA}: $r = 8$, $\alpha = 16$, $\text{lr} = 1 \times 10^{-4}$, replay\_ratio~=~0,2 (доля примеров из исходного домена в каждом batch для предотвращения catastrophic forgetting).


\subsection{Результаты и анализ}

\textbf{Фазы~1--2: Devign.}
Воспроизведение ReGVD baseline подтвердило корректность реализации (Accuracy~=~0,646, F1~=~0,571, ROC-AUC~=~0,704). Введение LoRA-адаптеров, embedding head и FAISS-индекса не привело к значимому изменению качества: отклонение по всем метрикам составило менее 0,5~п.п. (табл.~\ref{tab:summary}, столбцы Ф.~1 и Ф.~2). Данный результат подтверждает совместимость добавленных компонентов с базовой архитектурой: контрастивная компонента $\mathcal{L}_{\text{SupCon}}$ формирует пространство эмбеддингов, ортогональное задаче бинарной классификации, не внося деструктивных интерференций. Незначительное увеличение числа обучаемых параметров (с ${\sim}12{,}7$M до ${\sim}13{,}0$M) не приводит к переобучению благодаря регуляризирующему эффекту LoRA.

\textbf{Фазы~3--4: MegaVul (domain transfer).}
Zero-shot перенос модели, обученной на сбалансированном Devign, на несбалансированный MegaVul показал существенное падение метрик (ROC-AUC: $0{,}700 \to 0{,}572$, F1: $0{,}575 \to 0{,}124$), подтверждая наличие значительного domain gap между датасетами (табл.~\ref{tab:summary}, столбцы Ф.~3 и Ф.~4). Основные причины падения: (1)~различие в распределении классов (45,6\% vs. 5,1\% уязвимых), (2)~различие в стиле кода (4~проекта vs. 992~проекта), (3)~различие в типах уязвимостей (Devign не содержит CWE-аннотаций, MegaVul~--- 169~типов CWE).

LoRA-адаптация с обучением лишь 0,12\% параметров (${\sim}152{,}8$K из ${\sim}125$M) существенно восстановила качество: ROC-AUC~$+0{,}179$, Precision увеличился в~5~раз (с 0,073 до 0,367), F1~--- в~2,3~раза (с 0,124 до 0,284). Accuracy выросла с 0,621 до 0,928, что объясняется корректной калибровкой модели на несбалансированном распределении: модель научилась предсказывать отрицательный класс для большинства функций, что соответствует реальному соотношению классов.

Снижение Recall ($0{,}428 \to 0{,}231$)~--- ожидаемый результат калибровки: при ${\sim}5{,}1$\% уязвимых функций высокий Recall неизбежно сопровождается низкой Precision. Выбранный компромисс (Precision~=~0,367, Recall~=~0,231) соответствует практической задаче, где стоимость ложного срабатывания (ненужная sandbox-верификация) существенно ниже стоимости пропущенной уязвимости, но чрезмерно высокий false positive rate делает систему непрактичной.

\textbf{Сводные результаты.}
Табл.~\ref{tab:summary} обобщает результаты всех четырёх экспериментальных фаз. Ключевые наблюдения: (1)~гибридная архитектура нейтральна к базовому качеству на исходном домене; (2)~domain gap между сбалансированным и несбалансированным датасетами является значительным; (3)~LoRA-адаптация эффективно преодолевает domain gap при минимальных вычислительных затратах.

\begin{table}
\centering
\caption{Сводные результаты экспериментальных фаз}
\label{tab:summary}
\small
\begin{tabular}{|l|*{4}{c|}}
\hline
\textbf{Метрика} & \textbf{Ф.~1: ReGVD} & \textbf{Ф.~2: Hybrid} & \textbf{Ф.~3: Zero-shot} & \textbf{Ф.~4: + LoRA} \\
& \footnotesize (Devign) & \footnotesize (Devign) & \footnotesize (MegaVul) & \footnotesize (MegaVul) \\
\hline
Accuracy & 0,646 & 0,645 & 0,621 & \textbf{0,928} \\
ROC-AUC & 0,704 & 0,700 & 0,572 & \textbf{0,751} \\
F1 & 0,571 & 0,575 & 0,124 & \textbf{0,284} \\
Precision & 0,603 & 0,599 & 0,073 & \textbf{0,367} \\
Recall & 0,543 & 0,552 & 0,428 & 0,231 \\
Обуч. парам. & ${\sim}12{,}7$M & ${\sim}13{,}0$M & --- & ${\sim}$\textbf{152,8K} \\
\hline
\end{tabular}
\end{table}

\textbf{Системная производительность.}
Табл.~\ref{tab:performance} содержит характеристики системной производительности. Прирост латентности инференса составил ${\sim}21$\% (с 28~мс до 34~мс) за счёт embedding head и FAISS-поиска. Данное увеличение приемлемо для CI/CD-интеграции: при throughput ${\sim}29$~функций/с полное сканирование проекта из 10\,000~функций занимает менее 6~минут. VRAM при инференсе вырос незначительно (с 1,9~ГБ до 2,1~ГБ), что позволяет выполнять инференс даже на бюджетных GPU.

Основной ресурсный выигрыш достигается при дообучении: VRAM снижается с ${\sim}14{,}2$~ГБ до ${\sim}7{,}8$~ГБ (в~1,8~раза), энергозатраты~--- с ${\sim}1{,}92$~кВт$\cdot$ч до ${\sim}0{,}23$~кВт$\cdot$ч (в~${\sim}8$~раз). В пересчёте на эквивалент CO$_2$ это составляет снижение с ${\sim}920$~г до ${\sim}110$~г на цикл дообучения.

\begin{table}
\centering
\caption{Системная производительность (NVIDIA RTX 3090)}
\label{tab:performance}
\small
\begin{tabular}{|l|c|c|}
\hline
\textbf{Показатель} & \textbf{ReGVD baseline} & \textbf{Гибридный} \\
\hline
Inference latency, мс & ${\sim}28$ & ${\sim}34$ \\
Throughput, функций/с & ${\sim}36$ & ${\sim}29$ \\
VRAM при инференсе, ГБ & ${\sim}1{,}9$ & ${\sim}2{,}1$ \\
VRAM при дообучении, ГБ & ${\sim}14{,}2$ & ${\sim}$\textbf{7,8} \\
Hot-update (100 эмб.), с & --- & ${\sim}0{,}18$ \\
Энергозатраты дообучения, кВт$\cdot$ч & ${\sim}1{,}92$ & ${\sim}$\textbf{0,23} \\
Эквивалент CO$_2$, г & ${\sim}920$ & ${\sim}$\textbf{110} \\
\hline
\end{tabular}
\end{table}

\textbf{Покрытие типов уязвимостей.}
Табл.~\ref{tab:vuln_coverage} демонстрирует матрицу покрытия: модуль sandbox-верификации поддерживает 18~типов уязвимостей (23~записи CWE). Для каждого типа указаны наличие exploit-шаблонов, generic-генераторов, используемый профиль санитайзера и поддержка taint-анализа. Наиболее полное покрытие имеют уязвимости работы с памятью (buffer overflow, UAF, format string): для них доступны специализированные шаблоны, generic-генераторы, taint-анализ и санитайзер ASan. Уязвимости concurrency (data race, deadlock) покрыты через TSan, но без taint-анализа, поскольку taint-анализ ориентирован на потоки данных, а не на межпотоковые взаимодействия.

\begin{table}
\centering
\caption{Матрица покрытия типов уязвимостей}
\label{tab:vuln_coverage}
\small
\begin{tabular}{|l|l|c|c|c|c|}
\hline
\textbf{Тип уязвимости} & \textbf{CWE} & \textbf{Шабл.} & \textbf{Generic} & \textbf{Санит.} & \textbf{Taint} \\
\hline
Buffer Overflow (stack) & CWE-121 & 3 & \checkmark & ASan & \checkmark \\
Buffer Overflow (heap) & CWE-122 & 2 & \checkmark & ASan & \checkmark \\
Use-After-Free & CWE-416 & 2 & \checkmark & ASan & \checkmark \\
Double-Free & CWE-415 & --- & \checkmark & ASan & \checkmark \\
Format String & CWE-134 & 2 & \checkmark & ASan & \checkmark \\
Integer Overflow & CWE-190 & 2 & \checkmark & UBSan & --- \\
Null Pointer Deref & CWE-476 & 2 & \checkmark & ASan & --- \\
Off-by-One & CWE-193 & 2 & \checkmark & ASan & \checkmark \\
Division by Zero & CWE-369 & 2 & \checkmark & UBSan & --- \\
Uninit. Read & CWE-457 & 2 & \checkmark & MSan & --- \\
Data Race & CWE-362 & 1 & \checkmark & TSan & --- \\
Deadlock & CWE-833 & --- & --- & TSan & --- \\
\hline
\end{tabular}
\end{table}

\textbf{Результаты LLM-триажа.}
LLM-голова триажа (Claude Sonnet 4.6 + CodeQL) оценена на отложенной выборке из 848~примеров (табл.~\ref{tab:llm_triage}). Accuracy составила 69,6\%, при этом полнота на классе уязвимостей~--- 81,6\% (346 из 424), что соответствует консервативной калибровке: стоимость пропуска реальной уязвимости (false negative) существенно выше стоимости ложного срабатывания (false positive). Полнота на классе безопасных функций составила 57,1\% (242 из 424), что указывает на тенденцию LLM к over-reporting уязвимостей.

Доля неопределённых ответов (uncertain) составила лишь 0,35\% (3~из~848), что подтверждает стабильность формата ответов и пригодность LLM-головы для автоматизированной интеграции без ручного постпроцессинга. Стоимость одного запроса составляет ${\sim}\$0{,}006$, что делает LLM-триаж экономически целесообразным даже при масштабном сканировании.

\begin{table}
\centering
\caption{Результаты LLM-головы триажа ($N = 848$)}
\label{tab:llm_triage}
\small
\begin{tabular}{|l|c|c|}
\hline
\textbf{Показатель} & \textbf{Значение} & \textbf{Комментарий} \\
\hline
Accuracy & 69,6\% & без uncertain \\
True positives & 346/424 & Полнота: 81,6\% \\
True negatives & 242/424 & Полнота: 57,1\% \\
Uncertain & 3 & 0,35\% \\
Стоимость запроса & \$0,006 & Claude Sonnet 4.6 \\
\hline
\end{tabular}
\end{table}


% ====================================================================
%              6. ЗАКЛЮЧЕНИЕ
% ====================================================================
\section{Заключение}

В работе предложена интегрированная архитектура сканера уязвимостей GNN-ReGVD+, объединяющая четыре модуля: GNN-детекцию на основе модели ReGVD с LoRA-адаптацией и гибридным скорингом через FAISS, извлечение exploit-шаблонов из векторной базы, каскадную LLM-адаптацию шаблонов с итеративным уточнением и мульти-санитайзерную sandbox-верификацию (ASan+UBSan, MSan, TSan) с coverage-guided фаззингом через LibFuzzer. Параллельно функционирует LLM-голова триажа на основе Claude Sonnet 4.6 с графовым контекстом CodeQL.

Экспериментальная оценка на датасетах Devign (27\,318~функций) и MegaVul (339\,548~функций) подтвердила четыре ключевых результата: (а)~гибридная архитектура не снижает базовую производительность модели на исходном домене ($\Delta < 0{,}5$~п.п. по всем метрикам); (б)~LoRA-адаптация 0,12\% параметров (${\sim}152{,}8$K из ${\sim}125$M) обеспечивает прирост ROC-AUC на $+0{,}179$ при domain transfer со снижением VRAM в~1,8~раза и энергозатрат в~8~раз; (в)~мульти-санитайзерная верификация покрывает 18~типов уязвимостей (23~записи CWE) с четырёхуровневой классификацией доказательств; (г)~LLM-голова триажа обеспечивает полноту 81,6\% на классе уязвимостей при стоимости \$0,006 за запрос. Реализация (9\,372~строки кода Python, 203~автоматических теста, CLI-интерфейс) ориентирована на интеграцию в CI/CD-инфраструктуру.

К ограничениям системы относятся: поддержка только C/C++, невысокий Recall на несбалансированных данных (0,231 на MegaVul), Docker overhead (${\sim}2$--5~с на harness), Accuracy LLM-триажа 69,6\% и ограниченная база из 20~exploit-шаблонов. Перспективными направлениями развития являются расширение языковой поддержки (Python, Java), применение стратегий борьбы с дисбалансом классов (focal loss, class-balanced sampling), снижение sandbox overhead через container pool, интеграция символического выполнения (KLEE) и расширение базы шаблонов из CVEfixes и DiverseVul.


%
\begin{credits}
\subsubsection{\discintname}
Авторы заявляют об отсутствии конфликта интересов.
\end{credits}
%
% ---- Bibliography ----
%
\begin{thebibliography}{17}

\bibitem{nguyen2022regvd}
Nguyen~V.-A., Nguyen~D.\,Q., Nguyen~V., Le~T., Tran~Q.\,H., Phung~D.:
ReGVD: Revisiting Graph Neural Networks for Vulnerability Detection.
In: Proc. ICSE '22 Companion, pp.~164--165 (2022)

\bibitem{zhou2019devign}
Zhou~Y., Liu~S., Siow~J., Du~X., Liu~Y.:
Devign: Effective Vulnerability Identification by Learning Comprehensive Program Semantics via Graph Neural Networks.
arXiv preprint arXiv:1909.03496 (2019)

\bibitem{guo2021graphcodebert}
Guo~D., Ren~S., Lu~S. et~al.:
GraphCodeBERT: Pre-training Code Representations with Data Flow.
arXiv preprint arXiv:2009.08366 (2021)

\bibitem{hu2021lora}
Hu~E.\,J., Shen~Y., Wallis~P. et~al.:
LoRA: Low-Rank Adaptation of Large Language Models.
arXiv preprint arXiv:2106.09685 (2021)

\bibitem{khosla2020supcon}
Khosla~P. et~al.:
Supervised Contrastive Learning.
In: NeurIPS, vol.~33, pp.~18661--18673 (2020)

\bibitem{johnson2019faiss}
Johnson~J., Douze~M., J\'{e}gou~H.:
Billion-scale Similarity Search with GPUs.
IEEE Trans. Big Data \textbf{7}(3), 535--547 (2019)

\bibitem{ding2024}
Ding~Y., Fu~Y., Ibrahim~O. et~al.:
Vulnerability Detection with Code Language Models: How Far Are We?
arXiv preprint arXiv:2403.18624 (2024)

\bibitem{sultana2024}
Sultana~Sh., Afrin~S., Eisti~N.\,Y.:
Code Vulnerability Detection: A Comparative Analysis of Emerging Large Language Models.
arXiv preprint arXiv:2409.10490 (2024)

\bibitem{makhyari2024}
Makhyari~A.\,A.:
Harnessing the Power of LLMs in Source Code Vulnerability Detection.
arXiv preprint arXiv:2408.03489 (2024)

\bibitem{li2025neurosymbolic}
Li~P. et~al.:
Neuro-symbolic Static Analysis with LLM-generated Vulnerability Patterns.
arXiv preprint arXiv:2504.16057 (2025)

\bibitem{gajjar2026vbf}
Gajjar~M.:
Verify Before You Fix: Execution-Based Validation for LLM-Generated Vulnerability Patches.
arXiv preprint (2026)

\bibitem{agrawal2025sast}
Agrawal~A., Ahi~P.:
SAST-Genius: LLM-Based Triage and PoC Generation for SAST Alerts.
arXiv preprint (2025)

\bibitem{ni2024megavul}
Ni~C., Shen~L., Yang~X., Zhu~Y., Wang~S.:
MegaVul: A C/C++ Vulnerability Dataset with Comprehensive Code Representations.
In: Proc. MSR~'24, pp.~738--742 (2024)

\bibitem{chen2024uaf}
Chen~X., Li~Y., Wang~Z.:
Use-After-Free Vulnerability Detection via Deep Graph Learning.
IEEE Trans. Software Engineering \textbf{50}(4), 1123--1139 (2024)

\bibitem{nist800}
NIST:
Risk Management Guide for Information Technology Systems.
NIST SP 800-30 Rev.~1 (2023)

\bibitem{owasp2024}
OWASP Foundation:
OWASP Top 10:2024~--- The Ten Most Critical Web Application Security Risks (2024)

\bibitem{seacord2021}
Seacord~R.\,C.:
Secure Coding in C and C++, 3rd edn. Addison-Wesley, Boston (2021)

\end{thebibliography}

\end{document}
